\section*{Capítulo \ref{ch:g11:seismic-structure} --- As ondas sísmicas revelam o interior}

\begin{solution}{exo:g11:seismic-structure:1}
Uma \omterm{def:g11:seismic-structure:body-waves}{onda P} comprime e estica a rocha, que se move ao longo da direção de
propagação; uma \omterm{def:g11:seismic-structure:body-waves}{onda S} a cisalha, e a rocha se move perpendicularmente a essa
direção. As \omterm{def:g11:seismic-structure:body-waves}{ondas P} atravessam sólidos e líquidos; as \omterm{def:g11:seismic-structure:body-waves}{ondas S}, só sólidos. A
\omterm{def:g11:seismic-structure:body-waves}{onda P} é registrada primeiro, porque é a mais rápida das duas.
\end{solution}

\begin{solution}{exo:g11:seismic-structure:2}
$40 \times 111.2 = \qty{4448}{km}$, cerca de \qty{4450}{km}. $8900/111.2
\approx \ang{80}$.
\end{solution}

\begin{solution}{exo:g11:seismic-structure:3}
\qty{456}{s} são 7~minutos e 36~segundos; \qty{822}{s} são 13~minutos e
42~segundos. O atraso S--P é de $822 - 456 = \qty{366}{s}$, isto é,
6~minutos e 6~segundos.
\end{solution}

\begin{solution}{exo:g11:seismic-structure:4}
$\sin i_c = 6.0/8.0 = 0.75$, logo $i_c \approx \ang{48.6}$.
\end{solution}

\begin{solution}{exo:g11:seismic-structure:5}
O núcleo externo tem $5150 - 2891 = \qty{2259}{km}$ de espessura. O raio do
núcleo é $6371 - 2891 = \qty{3480}{km}$, e $3480/6371 \approx 0.55$: um pouco
mais da metade do raio da Terra.
\end{solution}

\begin{solution}{exo:g11:seismic-structure:6}
6~minutos e 6~segundos são \qty{366}{s}: o terremoto está a \ang{40}, isto
é, $40 \times 111.2 \approx \qty{4450}{km}$. Um único atraso só dá a
distância: o epicentro pode estar em qualquer ponto de um círculo desse raio
em torno da estação. São precisas mais duas estações para encontrar o único
ponto que fica à distância certa das três.
\end{solution}

\begin{solution}{exo:g11:seismic-structure:7}
$i_c \approx \ang{48.6}$ e $\cos i_c \approx 0.661$. Onda direta: $200/6.0
\approx \qty{33.3}{s}$. \omterm{ex:g11:seismic-structure:sp}{Onda frontal}: $200/8.0 + 2 \times 35 \times 0.661/6.0
\approx 25.0 + 7.7 = \qty{32.7}{s}$. A \omterm{ex:g11:seismic-structure:sp}{onda frontal} chega primeiro, cerca de
\qty{0.6}{s} antes.
\end{solution}

\begin{solution}{exo:g11:seismic-structure:8}
$\tan i_c \approx 1.134$: a distância crítica é $2 \times 35 \times 1.134
\approx \qty{79}{km}$. A \omterm{prop:g11:seismic-structure:crossover}{distância de cruzamento} é $2 \times 35 \times
\sqrt{14.0/2.0} = 70\sqrt 7 \approx \qty{185}{km}$. A estação do
\cref{exo:g11:seismic-structure:7} fica a \qty{200}{km}, além do
cruzamento: a \omterm{ex:g11:seismic-structure:sp}{onda frontal} chega primeiro, como se encontrou.
\end{solution}

\begin{solution}{exo:g11:seismic-structure:9}
$\sin i_c = 6.5/8.0 = 0.8125$, $i_c \approx \ang{54.3}$ e $\cos i_c \approx
0.583$. \omterm{prop:g11:seismic-structure:head-wave}{Tempo de intercepto}: $2 \times 6 \times 0.583/6.5 \approx \qty{1.1}{s}$.
\omterm{prop:g11:seismic-structure:crossover}{Distância de cruzamento}: $12\sqrt{14.5/1.5} \approx \qty{37}{km}$. Sob um
continente, o cruzamento fica perto de \qty{165}{km}, de modo que estações
a algumas dezenas de quilômetros umas das outras captam as duas retas; sob o
oceano, a \omterm{ex:g11:seismic-structure:sp}{onda frontal} já chega primeiro além de cerca de \qty{37}{km}, e a
reta da onda direta tem apenas algumas dezenas de quilômetros de
comprimento: as estações precisam estar a alguns quilômetros umas das outras
para traçá-la.
\end{solution}

\begin{solution}{exo:g11:seismic-structure:10}
$\sin r = (8.06/13.7)\sin\ang{30} \approx 0.294$, logo $r \approx \ang{17.1}$:
o raio se inclina em direção à perpendicular, isto é, em direção ao centro
da Terra. Ele atravessa, portanto, o núcleo por um longo trajeto perto do seu
meio, e volta a sair bem longe, do outro lado: os raios P que entram no
núcleo saem muito mais longe do terremoto que os raios que passam logo ao
lado dele.
\end{solution}

\begin{solution}{exo:g11:seismic-structure:11}
$12\,742/13.7 \approx \qty{930}{s}$, cerca de 15.5~minutos, enquanto a onda
leva \qty{1210}{s}, cerca de 20~minutos: ela chega uns 280~segundos
atrasada. Mesmo na velocidade mais alta do manto, a onda chegaria muito
antes, de modo que o meio da Terra precisa freá-la. Oldham concluiu que a
Terra tem um núcleo central no qual as \omterm{def:g11:seismic-structure:body-waves}{ondas P} se propagam mais devagar que
no manto.
\end{solution}

\begin{solution}{exo:g11:seismic-structure:12}
(1) Pela proposição, $t_i = 2h\cos i_c/v_1$, e $\cos i_c = \sqrt{1 -
v_1^2/v_2^2} = \sqrt{v_2^2 - v_1^2}/v_2$. Logo
\[
  t_i = \frac{2h\sqrt{v_2^2 - v_1^2}}{v_1 v_2},
  \qquad\text{portanto}\qquad
  h = \frac{t_i\,v_1 v_2}{2\sqrt{v_2^2 - v_1^2}}.
\]
(2) $h = 6.23 \times 6.2 \times 8.1/\bigl(2\sqrt{65.61 -
38.44}\bigr) \approx 312.9/10.43 \approx \qty{30.0}{km}$.
\end{solution}

\begin{solution}{exo:g11:seismic-structure:13}
(1) $(3480/6371)^3 \approx 0.163$: o núcleo é cerca de \qty{16}{\percent} do
volume. (2) A densidade média é a média ponderada pelos volumes: $5.52 = 0.163
\rho_c + 0.837 \times 4.45$, logo $\rho_c \approx (5.52 - 3.72)/0.163 \approx
\qty{11.0}{g/cm^3}$. (3) Isso é mais que \qty{9.9}{g/cm^3}, como deve ser,
já que a densidade do núcleo aumenta em direção ao centro (até
\qty{13.1}{g/cm^3}); é cerca de quatro vezes a densidade da crosta,
\qty{2.6}{g/cm^3}. O núcleo é feito de um material muito mais denso que as
rochas da superfície.
\end{solution}

\begin{solution}{exo:g11:seismic-structure:14}
(1) $105 \times 111.2 \approx \qty{11700}{km}$ e $143 \times 111.2 \approx
\qty{15900}{km}$. (2) $(\cos\ang{105} - \cos\ang{143})/2 \approx (-0.259 +
0.799)/2 \approx 0.27$: cerca de \qty{27}{\percent} da superfície, um anel
em torno do ponto oposto ao terremoto. (3) A \omterm{def:g11:seismic-structure:shadow-zone}{zona de sombra} é uma faixa de
distâncias ao epicentro: ela se desloca com o terremoto. Cada terremoto tem
o seu próprio anel.
\end{solution}

\begin{solution}{exo:g11:seismic-structure:15}
(1) $h = 9.5 \times 6.3 \times 8.1/\bigl(2\sqrt{65.61 - 39.69}\bigr) \approx
484.8/10.18 \approx \qty{48}{km}$. (2) $\sin i_c = 6.3/8.1 \approx 0.778$,
$\cos i_c \approx 0.629$, e $t_i = 2 \times 30 \times 0.629/6.3 \approx
\qty{6.0}{s}$. (3) A crosta sob a cordilheira tem cerca de \qty{48}{km} de
espessura, mais que a média continental de cerca de \qty{41}{km} e perto da
média das cadeias de montanhas: as cordilheiras altas se apoiam numa crosta
espessa.
\end{solution}

\begin{solution}{pb:g11:seismic-structure:1}
\textbf{1.} Os pontos Pg ficam numa reta; os pontos Pn ficam em outra reta,
menos inclinada.

\textbf{2.} O terremoto está a uma profundidade muito rasa sob o
epicentro: uma estação à distância zero registraria a onda imediatamente.

\textbf{3.} $(32.3 - 3.3)/(200 - 20) = 29.0/180 \approx \num{0.161}$~\unit{s}
por quilômetro.

\textbf{4.} $v_1 = 180/29.0 \approx \qty{6.2}{km/s}$.

\textbf{5.} $3.3 + 280 \times 0.161 \approx \qty{48.4}{s}$.

\textbf{6.} $(43.2 - 18.6)/(300 - 100) = 24.6/200 = \num{0.123}$~\unit{s}
por quilômetro, logo $v_2 = 200/24.6 \approx \qty{8.13}{km/s}$.

\textbf{7.} Ela é próxima de \qty{8.0}{km/s}: sob a crosta fica a rocha do
topo do manto.

\textbf{8.} $\sin i_c = 6.21/8.13 \approx 0.763$, logo $i_c \approx
\ang{49.8}$.

\textbf{9.} A \omterm{ex:g11:seismic-structure:sp}{onda frontal} desce com o ângulo $i_c$ e volta a subir com o
mesmo ângulo. O primeiro ponto em que ela pode voltar à superfície fica
$h\tan i_c$ além do ponto em que ela chegou ao \omterm{def:g11:seismic-structure:moho}{Moho}, que por sua vez fica a
$h\tan i_c$ da fonte: nenhuma estação mais próxima que $2h\tan i_c$ pode
recebê-la.

\textbf{10.} Um sismômetro registra todas as ondas que chegam até ele, não
só a primeira; a onda Pn chega depois da onda Pg, mas ainda aparece como uma
onda separada no registro.

\textbf{11.} A \qty{160}{km}, Pg (\qty{25.8}{s}) vem antes de Pn
(\qty{26.0}{s}); a \qty{180}{km}, Pn (\qty{28.5}{s}) vem antes de Pg
(\qty{29.0}{s}). A mudança acontece entre essas duas estações.

\textbf{12.} Pg: $t = 3.3 + 0.1611\,(x - 20) \approx 0.08 + 0.1611\,x$. Pn:
$t = 18.6 + 0.123\,(x - 100) = 6.30 + 0.123\,x$ ($x$ em quilômetros, $t$ em
segundos).

\textbf{13.} $0.08 + 0.1611\,x = 6.30 + 0.123\,x$ dá $x \approx
6.22/0.0381 \approx \qty{163}{km}$.

\textbf{14.} Pn, depois de $6.30 + 0.123 \times 250 \approx \qty{37.1}{s}$.

\textbf{15.} $t_i \approx \qty{6.3}{s}$.

\textbf{16.} $\cos i_c \approx 0.646$, logo $h = 6.21 \times 6.3/(2 \times
0.646) \approx \qty{30.3}{km}$.

\textbf{17.} $x_c = 2 \times 30.3 \times \sqrt{14.34/1.92} \approx 60.6 \times
2.73 \approx \qty{165}{km}$, a menos de \qty{2}{km} dos \qty{163}{km} da
questão~13.

\textbf{18.} $\tan i_c \approx 1.18$: $2 \times 30.3 \times 1.18 \approx
\qty{72}{km}$. A estação mais próxima que registra Pn fica a \qty{80}{km}, e
a de \qty{60}{km} não registra: os dois resultados concordam.

\textbf{19.} Os tempos são lidos a \qty{0.1}{s}, e um erro de
\qty{0.1}{s} só em $t_i$ muda $h$ em cerca de \qty{0.5}{km}; as velocidades,
tiradas de inclinações, também são incertas, e o terremoto não estava
exatamente na superfície. A crosta também não é perfeitamente uniforme. Uma
precisão de um ou dois quilômetros é o máximo que os dados permitem.

\textbf{20.} \textbf{O \omterm{def:g11:seismic-structure:moho}{Moho} fica a cerca de \qty{30}{km} sob esta região}:
a crosta aqui é do lado fino para um continente, cuja crosta tem em média
cerca de \qty{41}{km}; \qty{30}{km} é típico das regiões onde a crosta foi
estirada.
\end{solution}
